\inicializarSeccion{Resultados Obtenidos}
\textbf{Gráfica del perfil 2D y el sólido de revolución 3D}
\begin{figure}[H]
    \centering
    \includegraphics[width=0.65\textwidth]{\nat{parte4/resultado1.png}}
    \caption{Gráfica del perfil 2D y el sólido de revolución 3D.}
\end{figure}

\textbf{Gráfica que muestre claramente el espesor del producto}
\begin{figure}[H]
    \centering
    \includegraphics[width=0.65\textwidth]{\nat{parte4/resultado2.png}}
    \caption{Gráfica que muestra el espesor del producto.}
\end{figure}

\textbf{Gráfica o visualización que muestre que el producto cuenta con fondo}
\begin{figure}[H]
    \centering
    \begin{minipage}[t]{0.48\textwidth}
        \centering
        \includegraphics[width=0.65\textwidth]{\nat{parte4/resultado3_1.png}}
        \caption{Vista transversal del portavasos.}
        \label{fig:resultado3_1}
    \end{minipage}\hspace{0.02\textwidth}%
    \begin{minipage}[t]{0.48\textwidth}
        \centering
        \includegraphics[width=0.65\textwidth]{\nat{parte4/resultado3_2.png}}
        \caption{Vista semi-inclinada del portavasos.}
        \label{fig:resultado3_2}
    \end{minipage}
\end{figure}


\textbf{Captura de la ventana de comandos que imprima: volumen, masa y costo final (todos con sus unidades)}
\begin{figure}[H]
    \centering
    \includegraphics[width=0.65\textwidth]{\nat{parte4/resultado4.png}}
    \caption{Captura de la ventana de comandos con el volumen, masa y costo final.}
\end{figure}